\documentclass[12pt,a4paper]{article}
\raggedbottom

\usepackage{ugmath}
\usepackage{float}
\usepackage{placeins}
\usepackage[labelfont=bf,labelsep=space]{caption}
\usepackage{pdfpages}
\usepackage{tikz-cd}
\usepackage{amsthm}
\renewcommand{\Tr}{\operatorname{Tr}}
\newcommand{\llaves}[1]{\{#1\}}
\newcommand{\parentesis}[1]{\left(#1\right)}
\newcommand{\corchetes}[1]{\left[#1\right]}
\newcommand{\llavesgrande}[1]{\big\{#1\big\}}
\newcommand{\llavesgrandes}[1]{\left\{#1\right\}}


\begin{document}
        
    \thispagestyle{empty}
    \includepdf[pages=1]{main.pdf}

    \begin{problema}
        Si
        \begin{equation*}
            \begin{vmatrix}
                -3a_{11} & -6a_{13} & -3a_{12}\\
                a_{21} & 2a_{23} & a_{22}\\
                5a_{31} & 10a_{33} & 5a_{32}
            \end{vmatrix}
            =-15.
        \end{equation*}
        ¿Cuánto vale
        \begin{equation*}
            \begin{vmatrix}
                a_{11} & a_{12} & a_{13}\\
                a_{21} & a_{22} & a_{23}\\
                a_{31} & a_{32} & a_{33}
            \end{vmatrix}
            ?
        \end{equation*}
    \end{problema}

    \begin{problema}
        Una matriz $A\in M_{n\times n}(K)$ se llama ortogonal si $A^{-1}=A^{t}$. Demuestra que para $A$ ortogonal entonces $\det(A)=\pm1$.
    \end{problema}

    \begin{problema}
        Sea $A\in M_{n\times n}(K)$ una matriz cuya suma de todos los elementos de cada columna es igual a cero, entonces demuestra que $\det(A)=0$.
    \end{problema}

    \begin{problema}
        Una matriz cuadrada $B\in M_{n\times n}$ se llama antisimetrica si $B^{t}=B^{-1}$. Demuestra que si $B\in M_{n\times n}$ es antisimétrica y $n$ par, entonces $\det(B)=0$.
    \end{problema}

    \begin{problema}
        Encuentra la adjunta clásica de
        \begin{equation*}
            A=
            \begin{pmatrix}
                a_{11} & a_{12}\\
                a_{21} & a_{22}
            \end{pmatrix}
            \qquad
            B=
            \begin{pmatrix}
                3 & 6 & 7\\
                0 & 4 & 8\\
                0 & 0 & 5
            \end{pmatrix}.
        \end{equation*}
    \end{problema}

    \begin{problema}
        Demuestra que para cualquier $A\in M_{n\times n}(K)$, $\det(\adj(A))=[\det(A)]^{n-1}$.
    \end{problema}

    \begin{problema}
        Demuestra que si $A\in M_{n\times n}(K)$ es triangular superior entonces $\adj(A)$ también es triangular superior, concluye de esta forma que $A^{-1}$ es triangular superior.
    \end{problema}

\end{document}